% ============================================
% Bain & Company ADVANTAGE Program Cover Letter
% ============================================

\documentclass[11pt,letterpaper]{letter}

\usepackage{geometry}
\usepackage{microtype}
\usepackage[hidelinks]{hyperref}

\geometry{left=25mm,right=25mm,top=25mm,bottom=25mm}

\signature{Decory Edwards}
\address{
Department of Economics \\
Johns Hopkins University \\
Baltimore, MD 21218 \\
\texttt{dedwar65@jh.edu}
}

\begin{document}

\begin{letter}{
Recruiting Committee \\
Bain \& Company \\
ADVANTAGE Program
}

\opening{Dear Members of the Recruiting Committee,}

I am writing to express my strong interest in Bain \& Company’s ADVANTAGE Program. I am a Ph.D.\ candidate in Economics at Johns Hopkins University, advised by Professors Christopher D.~Carroll, Nicholas W.~Papageorge, and Stelios Fourakis, and I expect to complete my degree in May 2026. My training emphasizes quantitative analysis, economic modeling, and data-driven problem solving, and I am eager to apply these skills in a collaborative consulting environment focused on real-world impact.

My research agenda centers on understanding heterogeneity, incentives, and risk in complex economic systems. In my job market paper, I estimate persistent heterogeneity in individual returns to wealth using micro data from the Survey of Consumer Finances and embed these estimates into a structural model to study how dispersion in outcomes shapes aggregate dynamics. The project required translating large, noisy datasets into clear empirical patterns, stress-testing assumptions, and communicating results in a way that informs decision making under uncertainty. These are skills that closely parallel the analytical work performed in consulting engagements.

In a related project using the Health and Retirement Study, I study how beliefs and interpersonal trust affect portfolio choices and realized returns. I document systematic, non-linear relationships between behavioral traits and economic outcomes, highlighting how differences in information and decision-making processes can meaningfully alter performance. This work reflects my broader interest in diagnosing root causes, identifying mechanisms, and generating actionable insights from data—an approach that aligns well with Bain’s emphasis on rigorous analysis tied to practical recommendations.

I am particularly drawn to the ADVANTAGE Program because it is designed for candidates with advanced analytical training who are motivated to translate technical expertise into client-facing impact. Bain’s generalist model, team-oriented culture, and focus on mentorship strongly appeal to me, as does the opportunity to work across industries and problem domains. I am comfortable working with complex datasets, building quantitative frameworks, and presenting findings clearly to non-technical audiences, and I thrive in collaborative environments that value structured thinking and clear communication.

I would be enthusiastic about beginning my consulting career at Bain \& Company through the ADVANTAGE Program and contributing my quantitative background, problem-solving skills, and commitment to excellence to client teams. Thank you very much for your time and consideration.

\closing{Sincerely,}

\end{letter}
\end{document}
